\section{Comparison of Instance-Based and Associative Retrieval}

Retrieval mechanisms allow machine learning systems to access previously stored information without modifying the underlying model parameters. 
In classification systems, retrieval is commonly performed by comparing an input representation against stored examples or feature representations. 
Two relevant approaches are instance-based retrieval methods, such as k-nearest neighbors (KNN), and associative retrieval methods based on modern Hopfield networks.

\subsection{Retrieval-Augmented Classification}

Retrieval-Augmented Classification (RAC) extends conventional classification by introducing an external memory containing previously observed samples, representations, or class information. 
Instead of relying solely on knowledge encoded within the classifier parameters, predictions can be informed by retrieved information from this external memory.

A common approach is k-nearest neighbors (KNN), where inputs are first transformed into a continuous representation space. 
Given a query representation $q$ and a memory containing stored representations:

\begin{equation}
    M = \{(x_i,y_i)\}_{i=1}^{N}
\end{equation}

the similarity between the query and stored representations is calculated. 
In modern machine learning applications, cosine similarity is commonly used due to its effectiveness for comparing normalized embedding vectors:

\begin{equation}
    s(q,x_i)=\frac{q^Tx_i}{||q||_2||x_i||_2}
\end{equation}

The retrieved neighbours are then used to determine the predicted class:

\begin{equation}
    \hat{y}=
    \arg\max_{c\in C}
    \sum_{i\in N_k(q)}
    \mathbb{1}(y_i=c)
\end{equation}

where $N_k(q)$ denotes the set of the $k$ most similar stored representations.

Although KNN provides a simple and effective retrieval mechanism, the retrieval process is limited to selecting existing samples based on pairwise similarity. 
The retrieved information is not transformed into a new memory representation, and the contribution of each memory depends only on its similarity to the query.

\subsection{Associative Retrieval with Modern Hopfield Networks}

Modern Hopfield networks provide an alternative retrieval mechanism by treating memory access as an associative operation over a continuous representation space. 
Instead of selecting a fixed number of nearest neighbours, the query interacts with the complete memory set to produce a weighted combination of stored memories.

Given a memory matrix:

\begin{equation}
    M=[m_1,m_2,...,m_N]^T
\end{equation}

and a query representation $q$, retrieval is performed using:

\begin{equation}
    \xi^{new}
    =
    M^T \operatorname{softmax}(\beta Mq)
\end{equation}

where $\beta$ controls the concentration of the retrieval distribution.

The interaction term $Mq$ is based on the dot product between the query and stored memories. 
Unlike cosine similarity, the dot product retains information about vector magnitude in addition to directional similarity:

\begin{equation}
    m_i^Tq = ||m_i||_2||q||_2\cos(\theta)
\end{equation}

Consequently, modern Hopfield retrieval considers not only how well a memory aligns with the query, but also the magnitude of the interacting representations. 
This allows the retrieval process to incorporate contextual information contained within the representation space rather than relying solely on normalized similarity.

Therefore, while KNN performs explicit nearest-neighbour retrieval based on similarity between individual representations, modern Hopfield networks perform associative retrieval by dynamically combining information from multiple memories. 
The retrieved representation is not restricted to an existing stored sample but emerges from the interaction between the query and the complete memory space.

\subsection{Summary}

Both KNN-based retrieval and modern Hopfield networks use external memory representations to avoid storing all information within model parameters. 
However, they differ fundamentally in how memories are accessed. 
KNN performs direct similarity-based retrieval of existing examples, typically using cosine similarity between embedding vectors. 
Modern Hopfield networks instead perform associative retrieval through continuous interactions between queries and stored memories, allowing retrieved information to emerge as a weighted combination of stored representations.